\chapter{Zusammenfassung und Auswertung}
\label{chap:abschluss}

Die vorliegende Arbeit untersucht die Frage, ob und mit welcher Güte sich Log-Nachrichten kontinuierlich und ohne manuellen Pflegeaufwand gruppieren lassen. Ausgangspunkt ist die in Kapitel~\ref{chap:einf} dargelegte Beobachtung, dass Log-Nachrichten moderner Software-Systeme keinem einheitlichen Format folgen und dass klassische Ansätze wie regelbasierte Regex-Filter einen hohen manuellen Pflegeaufwand erfordern. Gesucht wurden daher Verfahren, die eingehende Log-Nachrichten automatisch Gruppen -- Clustern -- zuordnen und diese Zuordnung bei neuen Nachrichten fortlaufend aktualisieren, ohne manuellen Eingriff zu erfordern.

Zur Umsetzung dieses Ziels wurden zwei grundlegende Verfahrensklassen betrachtet. Kapitel~\ref{chap:ml} behandelt algorithmische Ansätze des maschinellen Lernens: Der partitionierende Algorithmus K-Means sowie der dichtebasierte Algorithmus DBSCAN operieren auf TF-IDF-Vektoren (Term Frequency-Inverse Document Frequency) und erfordern eine vollständige Neuverarbeitung bei jedem Durchlauf. Ergänzend werden mit Drain3 und Brain zwei spezialisierte Online-Parser betrachtet, die Log-Nachrichten iterativ einem wachsenden Template-Baum zuordnen und somit nativ für den kontinuierlichen Betrieb ausgelegt sind. Kapitel~\ref{chap:vek} stellt den vektorbasierten Ansatz vor: Jede Log-Nachricht wird mittels des vortrainierten Sentence-Transformer-Modells \texttt{all-MiniLM-L6-v2} in einen 384-dimensionalen semantischen Vektor überführt und in einer PostgreSQL-Datenbank mit \texttt{pgvector}-Erweiterung persistiert. Ein zweistufiger Clustering-Algorithmus ordnet neu eingehende Nachrichten zunächst bestehenden Clustern zu (Re-Clustering) und legt nur bei semantischer Neuartigkeit einen neuen Cluster an. Als weiterer Vergleichskandidat wird MinHash LSH (Locality-Sensitive Hashing) untersucht, das auf Signaturen statt auf Vektoren operiert.

Die praktische Realisierung aller Verfahren erfolgte in einer reproduzierbaren Docker-Umgebung bestehend aus einem Jupyter-Notebook-Container und einer PostgreSQL-Instanz mit \texttt{pgvector}-Erweiterung, wie in Kapitel~\ref{chap:real} beschrieben. Als Datenbasis dienen die drei Systemlog-Quellen \texttt{Linux}, \texttt{Apache} und \texttt{OpenSSH} aus dem Loghub-2k-Benchmark \cite{zhu2023loghublargecollectionlog}, die zu einem gemeinsamen Datensatz mit $10\,000$ annotierten Log-Nachrichten und $44$ Ground-Truth-Klassen zusammengeführt werden. Jeder Algorithmus wurde in fünf Durchläufen mit zufällig permutierter Eingangsreihenfolge evaluiert, um das Verhalten im kontinuierlichen Betrieb abzubilden. Die Clustering-Güte wird anhand von zwei komplementären Metriken beurteilt: dem Adjusted Rand Index~(ARI), der die Übereinstimmung mit der Ground Truth zufallskorrigiert misst, und dem V-Measure, das Homogenität und Vollständigkeit der gefundenen Cluster als harmonisches Mittel vereint.

Das folgende Kapitel fasst die Ergebnisse dieser Versuche vergleichend zusammen, diskutiert die Stärken und Schwächen der einzelnen Verfahren und leitet daraus Empfehlungen für den produktiven Einsatz ab.

\section{Vergleichende Auswertung der Clustering-Verfahren}
\label{sec:abschluss:auswertung}

Die in Kapitel~\ref{chap:eval} erhobenen Einzelergebnisse werden im Folgenden entlang der drei in Kapitel~\ref{chap:einf} definierten Bewertungskriterien -- Cluster-Qualität, Reihenfolge-Stabilität und strukturelle Eignung für den kontinuierlichen Betrieb -- vergleichend eingeordnet. Als Referenz dient Tabelle~\ref{tab:eval:gesamt}.

\paragraph{Gesamtranking nach ARI und V-Measure.}
Gemessen am ARI ergibt sich die Rangfolge Drain3 ($91{,}084\,\%$), pgvector ($79{,}273\,\%$), Brain ($73{,}280\,\%$), K-Means ($70{,}585\,\%$), MinHash~LSH ($43{,}378\,\%$) und DBSCAN ($26{,}085\,\%$). Beim V-Measure ergibt sich eine nahezu identische Reihenfolge: Drain3 führt mit $91{,}528\,\%$, gefolgt von pgvector ($88{,}451\,\%$), Brain ($88{,}102\,\%$), K-Means ($84{,}913\,\%$), MinHash~LSH ($64{,}890\,\%$) und DBSCAN ($59{,}905\,\%$). pgvector und Brain liegen im V-Measure mit weniger als einem halben Prozentpunkt nahezu gleichauf, fallen jedoch im ARI deutlich auseinander; das V-Measure bewertet Homogenität und Vollständigkeit symmetrisch, während der ARI über die paarweise Übereinstimmung empfindlicher auf die Cluster-Größenverteilung reagiert. DBSCAN bildet in beiden Metriken das Schlusslicht; die niedrigen Werte spiegeln eine vergleichsweise geringe Homogenität ($57{,}391\,\%$) bei ebenfalls niedriger Vollständigkeit ($62{,}649\,\%$) wider, die in der heterogenen Dichtestruktur des TF-IDF-Raums begründet liegen. Auf dem in dieser Arbeit verwendeten Loghub-Datensatz belegt der Online-Parser Drain3 in beiden Metriken den ersten Rang; die syntaktisch-strukturelle Template-Extraktion erweist sich hier als leistungsfähiger als die semantische Repräsentation über Sentence-Embeddings, die mit pgvector den zweiten Rang erreicht. Eine Übertragung dieser Rangfolge auf andere Log-Datensätze ist nicht ohne Weiteres zulässig (vgl.\ Abschnitt~\ref{sec:abschluss:limitierungen}).

\paragraph{Homogenität vs.\ Vollständigkeit: das zentrale Trade-off.}
Drei der sechs Verfahren erreichen eine Homogenität oberhalb von $90\,\%$: Brain ($98{,}731\,\%$), MinHash~LSH ($95{,}875\,\%$) und Drain3 ($91{,}851\,\%$). pgvector liegt mit $83{,}998\,\%$ knapp darunter, K-Means mit $80{,}762\,\%$. Die zugehörigen Vollständigkeitswerte unterscheiden sich jedoch erheblich -- von $49{,}041\,\%$ bei MinHash~LSH bis $93{,}433\,\%$ bei pgvector. Die Differenz zwischen Homogenität und Vollständigkeit beträgt bei MinHash~LSH $46{,}8$~Prozentpunkte und bei Brain $19{,}2$~Prozentpunkte; beide Verfahren erzeugen sehr viele, jeweils reine Cluster, die jedoch semantisch zusammengehörige Lognachrichten auf zahlreiche Partitionen aufteilen. Drain3 weist mit lediglich $0{,}6$~Prozentpunkten Differenz das ausgeglichenste Profil der in dieser Arbeit untersuchten Verfahren auf; vier der sechs Verfahren -- Drain3, pgvector, Brain und K-Means -- erreichen in beiden Teilmetriken Werte oberhalb von $70\,\%$. DBSCAN bildet den Sonderfall einer vergleichsweise niedrigen Homogenität ($57{,}391\,\%$) bei ebenfalls niedriger Vollständigkeit ($62{,}649\,\%$) -- ein Profil, das weder klar auf Über- noch auf Untersegmentierung hindeutet, sondern auf die heterogene Dichtestruktur des TF-IDF-Raums zurückzuführen ist. Die Ergebnisse zeigen, dass hohe Homogenität allein keine hinreichende Bedingung für eine gute Clustering-Güte darstellt; entscheidend ist das Gleichgewicht beider Teilmetriken.

\paragraph{Stabilität über Permutationsdurchläufe.}
Die Standardabweichung des ARI über die fünf Permutationsdurchläufe offenbart deutliche Unterschiede in der Reihenfolge-Stabilität. DBSCAN, Brain und MinHash~LSH weisen eine Standardabweichung nahe null auf -- bei den beiden erstgenannten aufgrund deterministischer Verarbeitung, bei DBSCAN aufgrund der Reihenfolge-Invarianz des Stapelverfahrens. Drain3 erreicht trotz seiner Online-Verarbeitung eine nahezu verschwindende Abweichung von $\pm\,0{,}120\,\%$; die positionsbasierte Verzweigung im Parse-Baum stabilisiert die Zuordnungen auch bei variierender Eingangsreihenfolge. pgvector zeigt mit $\pm\,5{,}293\,\%$ eine moderate Varianz, die aus dem inkrementellen Cluster-Aufbau resultiert: Die Reihenfolge, in der Lognachrichten eintreffen, beeinflusst die initialen Cluster-Schwerpunkte und damit nachfolgende Zuordnungsentscheidungen (vgl.\ Abschnitt~\ref{PG_Algorithmn}). K-Means weist mit $\pm\,8{,}369\,\%$ die höchste Varianz auf, bedingt durch die stochastische Zentroid-Initialisierung in Verbindung mit dem niedrigen Silhouette-Niveau im TF-IDF-Raum (vgl.\ Abschnitt~\ref{sec:eval:param:kmeans}). Für den produktiven Einsatz ist eine geringe Varianz wünschenswert, da wechselnde Ergebnisse bei identischem Datenbestand das Vertrauen in die Zuordnung untergraben.

\paragraph{Clusteranzahl und Ground-Truth-Nähe.}
Die erzeugten Clusteranzahlen spannen ein breites Spektrum auf: K-Means bildet mit $13$~Clustern die gröbste Partitionierung, gefolgt von Drain3 ($52$), pgvector ($80$), DBSCAN ($124$) und Brain ($151$). Bei $44$~Ground-Truth-Klassen im Datensatz erzeugt K-Means somit eine Untersegmentierung um den Faktor~$3{,}4$, während Brain um den Faktor~$3{,}4$ in die entgegengesetzte Richtung übersegmentiert. Drain3 liegt mit $52$~Clustern der Ground Truth am nächsten; die moderate Abweichung von acht zusätzlichen Clustern spiegelt sich in einem nahezu ausgeglichenen Profil mit hoher Homogenität ($91{,}851\,\%$) und ebenso hoher Vollständigkeit ($91{,}208\,\%$) wider. Die Ergebnisse bestätigen den erwarteten Zusammenhang: Verfahren mit deutlich zu wenigen Clustern erzielen hohe Vollständigkeit auf Kosten der Homogenität (K-Means), während Verfahren mit deutlich zu vielen Clustern die umgekehrte Tendenz zeigen (Brain, MinHash~LSH). pgvector und Drain3 treffen jeweils einen Kompromiss, der in beiden Teilmetriken oberhalb von $83\,\%$ liegt.

\paragraph{Eignung für kontinuierlichen Betrieb.}
Das in Kapitel~\ref{chap:einf} formulierte dritte Bewertungskriterium -- die Fähigkeit zur inkrementellen Verarbeitung mit persistierbarem Zustand -- trennt die untersuchten Verfahren in zwei Klassen. K-Means, DBSCAN und MinHash~LSH sind Stapelverfahren, die den gesamten Datensatz bei jeder Ausführung neu einlesen und verarbeiten müssen. Eine kontinuierliche Erweiterung um einzelne Lognachrichten ist ohne vollständige Neuberechnung nicht möglich; im produktiven Betrieb mit wachsendem Datenvolumen steigt der Rechenaufwand damit linear mit der Gesamtgröße des Log-Bestands. Drain3 und Brain verarbeiten Lognachrichten dagegen sequenziell und aktualisieren ihren Template-Baum inkrementell. Neue Nachrichten werden einem bestehenden Template zugeordnet oder erzeugen einen neuen Ast, ohne dass der bisherige Baumzustand neu berechnet werden muss; der Zustand kann zudem persistiert und über Neustarts hinweg wiederhergestellt werden. pgvector folgt einem vergleichbaren Prinzip: Jede neue Lognachricht wird in einen Embedding-Vektor überführt und gegen die in der PostgreSQL-Datenbank gespeicherten Cluster-Schwerpunkte abgeglichen; bei ausreichender Ähnlichkeit erfolgt die Zuordnung zu einem bestehenden Cluster, andernfalls wird ein neuer angelegt. Die Persistierung in der Datenbank ist systeminhärent. Von den sechs untersuchten Verfahren erfüllen somit nur Drain3, Brain und pgvector die Grundanforderung der Aufgabenstellung an den kontinuierlichen Betrieb.

% ----------------------------------------------------------------
\section{Beantwortung der Forschungsfrage}
\label{sec:abschluss:forschungsfrage}

Die in Kapitel~\ref{forschungs_frage} formulierte Forschungsfrage lautet, welches Verfahren sich für die automatisierte, kontinuierliche Gruppierung von Lognachrichten am besten eignet. Zur Beantwortung wurden drei Bewertungskriterien definiert: Cluster-Qualität (gemessen über ARI und V-Measure), Reihenfolge-Stabilität (gemessen über die Standardabweichung bei permutierter Eingangsreihenfolge) und strukturelle Eignung für den kontinuierlichen Betrieb (inkrementelle Verarbeitung mit persistierbarem Zustand). Im Folgenden wird die Forschungsfrage anhand dieser Kriterien beantwortet.

Von den sechs untersuchten Verfahren scheiden K-Means, DBSCAN und MinHash~LSH bereits am dritten Kriterium aus: Als Stapelverfahren erfordern sie eine vollständige Neuberechnung bei jeder hinzukommenden Lognachricht und verfügen über keinen persistierbaren Zustand für den inkrementellen Betrieb. Unabhängig von ihrer Cluster-Qualität sind sie damit für den in der Aufgabenstellung geforderten kontinuierlichen Einsatz strukturell nicht geeignet. Die Bewertung reduziert sich somit auf die drei verbleibenden Kandidaten: Drain3, Brain und pgvector.

Brain erzielt in der vorliegenden Evaluation zwar die höchste Homogenität aller untersuchten Verfahren ($98{,}731\,\%$), weist jedoch mit einem ARI von $73{,}280\,\%$ und einer Vollständigkeit von $79{,}540\,\%$ die niedrigste Gesamtgüte unter den drei Online-fähigen Verfahren auf. Die $151$~erzeugten Cluster übersteigen die $44$~Ground-Truth-Klassen um mehr als das Dreifache; die resultierende Übersegmentierung begrenzt den praktischen Nutzen, da ein nachgelagertes Zusammenführen der Cluster zusätzlichen Aufwand erfordert. Brain scheidet daher als bevorzugtes Verfahren aus.

Auf dem verwendeten Loghub-Datensatz fällt die Entscheidung zwischen Drain3 und pgvector zugunsten von Drain3 aus, das in allen drei Bewertungskriterien Vorteile aufweist. Drain3 erzielt mit einem ARI von $91{,}084\,\%$ und einem V-Measure von $91{,}528\,\%$ die höchsten Werte der hier untersuchten sechs Verfahren. Sowohl die Homogenität ($91{,}851\,\%$) als auch die Vollständigkeit ($91{,}208\,\%$) liegen oberhalb von $91\,\%$; mit einer Differenz von lediglich $0{,}6$~Prozentpunkten ist das Metrikprofil das ausgewogenste der untersuchten Verfahren. Die Standardabweichung beim ARI liegt mit $\pm\,0{,}120\,\%$ nahe null, sodass Drain3 trotz seines Online-Charakters nahezu deterministische Ergebnisse unabhängig von der Eingangsreihenfolge liefert -- eine Eigenschaft, die im produktiven Betrieb das Vertrauen in die Zuordnung stärkt. Die $52$~erzeugten Cluster liegen zudem mit nur acht Clustern Abweichung sehr nahe an den $44$~Ground-Truth-Klassen.

pgvector erreicht mit einem ARI von $79{,}273\,\%$ und einem V-Measure von $88{,}451\,\%$ die zweithöchsten Werte der untersuchten Verfahren und erzielt mit $93{,}433\,\%$ den höchsten Vollständigkeitswert dieser Evaluation. Die semantische Repräsentation über Sentence-Embeddings erkennt strukturell verwandte Lognachrichten auch dann zuverlässig, wenn sie sich in variablen Feldern wie IP-Adressen oder Zeitstempeln unterscheiden. Dem stehen jedoch eine geringere Homogenität ($83{,}998\,\%$) sowie eine Standardabweichung von $\pm\,5{,}293\,\%$ beim ARI gegenüber, die aus der Reihenfolgeabhängigkeit des inkrementellen Cluster-Aufbaus resultiert.

Auf Basis der drei Bewertungskriterien und bezogen auf den verwendeten Loghub-Datensatz ergibt sich folgende Antwort auf die Forschungsfrage: Der Online-Parser Drain3 stellt in dieser Evaluation das leistungsfähigste Verfahren dar. Er erzielt die höchste Cluster-Qualität, eine nahezu perfekte Reihenfolge-Stabilität und eignet sich strukturell für den kontinuierlichen Betrieb; zugleich erfordert er weder ein vortrainiertes Embedding-Modell noch eine Vektordatenbank, was den infrastrukturellen Aufwand minimiert. Der vektorbasierte Ansatz mit pgvector stellt eine starke Alternative dar, insbesondere wenn die Robustheit gegenüber stark variierenden Feldern (IP-Adressen, Zeitstempel, UUIDs) im Vordergrund steht -- hier nutzt pgvector die Stärke der semantischen Embeddings, die ein rein syntaktisch arbeitender Parser wie Drain3 prinzipiell nicht erreichen kann. Für den untersuchten Datensatz fällt die Empfehlung damit eindeutig auf Drain3; pgvector bleibt für Logs mit hoher syntaktischer Variabilität und einem ausgeprägten Bedarf an semantischer Ähnlichkeitserkennung eine sinnvolle Option. Die Übertragbarkeit beider Befunde auf andere Log-Quellen ist durch die in Abschnitt~\ref{sec:abschluss:limitierungen} dargelegten Einschränkungen begrenzt und wäre in einer weitergehenden Evaluation zu prüfen.

% ----------------------------------------------------------------
\section{Limitierungen}
\label{sec:abschluss:limitierungen}

\paragraph{Datensatz}
Die Evaluation basiert auf drei Quellen des Loghub-2k-Benchmarks~\cite{zhu2023loghublargecollectionlog} -- \texttt{Linux}, \texttt{Apache} und \texttt{OpenSSH} -- mit insgesamt $10\,000$ Lognachrichten und $44$~Ground-Truth-Klassen. Der Datensatz deckt damit ausschließlich klassische Systemlogs ab, die dem Syslog-Format nahestehen und ein vergleichsweise homogenes Vokabular aufweisen. Lognachrichten aus Projekt des Praktikums unterscheiden sich in Struktur, Länge und Variabilität von den verwendeten Quellen.

\paragraph{Embedding-Modell}
Der vektorbasierte Ansatz verwendet das vortrainierte Sentence-Transformer-Modell \texttt{all-MiniLM-L6-v2}, das auf allgemeinsprachlichen Text trainiert wurde (vgl.\ Abschnitt~\ref{sec:real:pgvector}). Lognachrichten weichen in Syntax und Semantik von natürlicher Sprache ab: Sie enthalten technische Bezeichner, hexadezimale Werte, Dateipfade und wiederkehrende Schlüsselwörter, die in allgemeinsprachlichen Text unterrepräsentiert sind. Ein domänenspezifisch feinjustiertes Modell -- etwa ein auf Log-Nachrichten nachtrainierter Transformer -- könnte die Trennschärfe der Embeddings insbesondere bei strukturell ähnlichen, aber semantisch unterschiedlichen Templates erhöhen. Die vorliegende Arbeit untersucht diesen Einfluss nicht; die berichteten Ergebnisse des pgvector-Verfahrens stellen damit eine untere Schranke für die mit Sentence-Embeddings erreichbare Clustering-Güte dar.

\paragraph{Skalierbarkeit.}
Alle Verfahren wurden auf einem Datensatz mit $10\,000$ Lognachrichten evaluiert. Produktive Log-Aggregatoren verarbeiten Datenvolumina, die um mehrere Größenordnungen darüber liegen -- typischerweise $10^6$ bis $10^8$ Nachrichten pro Tag. Für die Stapelverfahren K-Means und DBSCAN wächst der Rechenaufwand mit der Datensatzgröße; bei K-Means skaliert die Laufzeit linear mit der Anzahl der Datenpunkte und Cluster, bei DBSCAN dominiert die Nachbarschaftssuche mit einem Aufwand von $\mathcal{O}(n \log n)$ bei Verwendung räumlicher Indexstrukturen. Die Online-Verfahren Drain3 und Brain verarbeiten jede Nachricht in konstanter Zeit relativ zur Baumtiefe, wobei die Baumgröße mit der Anzahl unterschiedlicher Templates wächst. Beim pgvector-Ansatz bestimmen zwei Faktoren die Skalierbarkeit: die Inferenzzeit des Sentence-Transformer-Modells pro Lognachricht sowie die Suchzeit im HNSW-Index, die logarithmisch mit der Anzahl gespeicherter Vektoren skaliert~\cite{githubGitHubPgvectorpgvector}. Eine systematische Laufzeitanalyse, die das Verhalten der Verfahren bei wachsendem Datenvolumen quantifiziert, wurde im Rahmen dieser Arbeit nicht durchgeführt und bleibt als offene Fragestellung bestehen.

% ----------------------------------------------------------------
\section{Ausblick}
\label{sec:abschluss:ausblick}

\paragraph{Domänenspezifische Embeddings.}
% Finetuning von Sentence-Transformern auf Log-Korpora (z. B. LogBERT, LogRoberta)
% Erwartung: höhere Trennschärfe bei variablen Feldern

\paragraph{Hybride Verfahren.}
% Kombination von Drain3 (syntaktische Vorfilterung) + pgvector (semantische Nachverfeinerung)
% Drain3 als schneller First-Pass, pgvector für unsichere Zuordnungen

\paragraph{Skalierungstests.}
% Evaluation auf größeren Datensätzen (Loghub-2.0, industrielle Logs)
% Laufzeit- und Speicherbenchmarks für den kontinuierlichen Betrieb

\paragraph{Integration in Log-Management-Systeme.}
% Anbindung an bestehende Systeme (Elasticsearch, Grafana Loki)
% Echtzeit-Pipeline mit Kafka/Flink als Ingestion-Layer
